\documentclass[12pt,a4paper]{article}
\usepackage[utf8]{inputenc}
\usepackage[T1]{fontenc}
\usepackage{amsmath,amssymb,amsthm}
\usepackage{physics}
\usepackage{geometry}
\usepackage{enumitem}
\usepackage{hyperref}
\usepackage{fancyhdr}
\usepackage{titlesec}

\geometry{margin=1in}
\pagestyle{fancy}
\fancyhf{}
\rhead{Lecture 5}
\lhead{Advanced Quantum Mechanics}
\cfoot{\thepage}

\titleformat{\section}{\large\bfseries}{}{0em}{}
\titleformat{\subsection}{\normalsize\bfseries}{}{0em}{}

\newtheorem{definition}{Definition}

\title{\textbf{Lecture 5: Identical Particles and the Exclusion Principle} \\ \large Pauli Exclusion, Spin-Statistics, and Atomic Shell Structure}
\author{Based on Leonard Susskind's \textit{Advanced Quantum Mechanics}}
\date{}

\begin{document}
\maketitle
\thispagestyle{fancy}

% ============================================================
\section{1. Identical Particles in Quantum Mechanics}
% ============================================================

In classical mechanics, identical particles can always be distinguished by tracking their trajectories.  In quantum mechanics, this is impossible: there is no trajectory, and identical particles are truly \emph{indistinguishable}.

\begin{definition}[Identical particles]
Two particles are \textbf{identical} if no measurement can distinguish one from the other.  The physics must be invariant under interchange of the two particles.
\end{definition}

\subsection{1.1 Exchange symmetry}

Consider a two-particle state $\psi(x_1, x_2)$.  Exchanging particles 1 and 2:
\[
    P_{12}\,\psi(x_1, x_2) = \psi(x_2, x_1).
\]
Since $P_{12}^2 = I$, the eigenvalues of $P_{12}$ are $\pm 1$.  There are exactly two possibilities:
\begin{itemize}[nosep]
    \item \textbf{Symmetric}: $\psi(x_2, x_1) = +\psi(x_1, x_2)$ \quad (bosons).
    \item \textbf{Antisymmetric}: $\psi(x_2, x_1) = -\psi(x_1, x_2)$ \quad (fermions).
\end{itemize}

% ============================================================
\section{2. The Spin-Statistics Connection}
% ============================================================

One of the deepest results in physics:
\[
    \boxed{\text{Integer spin} \;\Longleftrightarrow\; \text{bosons (symmetric)}, \qquad \text{Half-integer spin} \;\Longleftrightarrow\; \text{fermions (antisymmetric)}.}
\]
This is not an assumption---it is a theorem in relativistic quantum field theory.  Key observation: a $2\pi$ rotation of a half-integer spin particle introduces a phase factor of $-1$, which is directly connected to the antisymmetry of the wave function.

% ============================================================
\section{3. The Pauli Exclusion Principle}
% ============================================================

\begin{definition}[Pauli exclusion principle]
No two identical fermions can occupy the same quantum state.
\end{definition}

\noindent\textbf{Proof from antisymmetry.}  If both particles are in state $\phi$:
\[
    \psi(x_1, x_2) = \phi(x_1)\phi(x_2).
\]
This is symmetric under exchange, so it cannot describe fermions.  The antisymmetrized version:
\[
    \psi_A(x_1, x_2) = \frac{1}{\sqrt{2}}\bigl[\phi(x_1)\phi(x_2) - \phi(x_2)\phi(x_1)\bigr] = 0.
\]
The state vanishes identically.  Two fermions \emph{cannot} be in the same state.

\subsection{3.1 Bosons: the opposite behavior}

Bosons have no such restriction.  In fact, bosons \emph{prefer} to congregate: the probability of finding two bosons in the same state is enhanced relative to distinguishable particles.  This is the origin of Bose--Einstein condensation, laser coherence, and the integer quantum Hall effect.

% ============================================================
\section{4. Multi-Particle States}
% ============================================================

\subsection{4.1 Two distinguishable particles}

If particle 1 can be in states $\{\phi_i\}$ and particle 2 in states $\{\chi_j\}$, a general two-particle state is:
\[
    \psi(x_1, x_2) = \sum_{i,j} c_{ij}\,\phi_i(x_1)\,\chi_j(x_2).
\]

\subsection{4.2 Two identical bosons}

The state must be symmetrized:
\[
    \psi_S(x_1, x_2) = \frac{1}{\sqrt{2}}\bigl[\phi_a(x_1)\phi_b(x_2) + \phi_a(x_2)\phi_b(x_1)\bigr].
\]

\subsection{4.3 Two identical fermions}

The state must be antisymmetrized:
\[
    \psi_A(x_1, x_2) = \frac{1}{\sqrt{2}}\bigl[\phi_a(x_1)\phi_b(x_2) - \phi_a(x_2)\phi_b(x_1)\bigr].
\]
This can be written as a \textbf{Slater determinant}:
\[
    \psi_A(x_1, x_2) = \frac{1}{\sqrt{2}}\begin{vmatrix} \phi_a(x_1) & \phi_b(x_1) \\ \phi_a(x_2) & \phi_b(x_2) \end{vmatrix}.
\]
The determinant automatically vanishes if $a = b$, enforcing the exclusion principle.

% ============================================================
\section{5. Atomic Shell Structure}
% ============================================================

\subsection{5.1 Filling order}

In a hydrogen-like atom, each electron state is labeled by quantum numbers $(n, \ell, m, m_s)$.  The Pauli exclusion principle dictates that each combination of quantum numbers can be occupied by at most one electron.

For a given $n$ and $\ell$:
\begin{itemize}[nosep]
    \item There are $2\ell + 1$ values of $m$.
    \item Each $m$ has 2 spin states ($m_s = \pm\frac{1}{2}$).
    \item Total: $2(2\ell + 1)$ electrons per subshell.
\end{itemize}

\subsection{5.2 Shell capacities}

\begin{center}
\begin{tabular}{c c c c}
\textbf{Subshell} & $\ell$ & \textbf{States} & \textbf{Max electrons} \\
\hline
$s$ & 0 & 1 & 2 \\
$p$ & 1 & 3 & 6 \\
$d$ & 2 & 5 & 10 \\
$f$ & 3 & 7 & 14 \\
\end{tabular}
\end{center}

\subsection{5.3 The periodic table}

The structure of the periodic table arises directly from the Pauli exclusion principle applied to atomic electrons.  As atomic number $Z$ increases, electrons fill successively higher energy levels.  The chemical properties of an element are determined primarily by the number of electrons in the outermost (valence) shell:
\begin{itemize}[nosep]
    \item Noble gases: filled shells, chemically inert.
    \item Alkali metals: one electron beyond a filled shell, highly reactive.
    \item Halogens: one electron short of a filled shell, highly reactive.
\end{itemize}

\medskip
\noindent\textit{The Pauli exclusion principle, combined with the shell structure of atoms, explains the periodicity of chemical properties---one of the great triumphs of quantum mechanics.  Next, we move to the formalism of second quantization, which provides a natural language for many-particle systems.}

\end{document}
